\documentclass[11pt,a4paper]{article}

\usepackage[utf8]{inputenc}
\usepackage[T1]{fontenc}
\usepackage{lmodern}
\usepackage[margin=2.5cm]{geometry}
\usepackage{parskip}
\usepackage{ragged2e}
\usepackage[hidelinks]{hyperref}

\pagestyle{empty}
\setlength{\parindent}{0pt}

\begin{document}

% Candidate information
{\Large\bfseries John Doe}\\[0.3cm]


\href{mailto:john@example.com}{john@example.com}\\



+33 6 00 00 00 00\\



Paris\\


\vspace{1cm}

% Recipient
\begin{flushright}

\textbf{Natixis}\\



Paris\\


\vspace{0.5cm}

13 novembre 2026

\end{flushright}

\vspace{0.8cm}

% Subject
\textbf{Subject: Candidature spontanée – Stage en Data Engineering}

\vspace{0.6cm}

% Salutation

Madame, Monsieur,


\vspace{0.5cm}

% Cover letter body
\justifying


Je vous adresse ma candidature spontanée pour un stage en Data Engineering au sein de Natixis, à Paris. Attiré par les enjeux de la donnée dans le secteur bancaire et par l’importance des infrastructures de données, je souhaite contribuer à vos équipes.

\par

Actuellement étudiant en Master en informatique (Université Exemple, 2024–2026), j’ai consolidé des bases en Python, SQL et Linux. Lors de mon stage en 2025 chez Entreprise Exemple en tant que stagiaire développement, j’ai réalisé le développement de scripts en Python et la rédaction de requêtes SQL.

\par

Je serais ravi d’échanger afin de vous présenter ma motivation et d’explorer la manière dont je pourrais soutenir vos activités data.

\par


\vspace{0.7cm}

% Closing

Cordialement,


\vspace{0.5cm}

\textbf{John Doe}

\end{document}

